% --- CORRESPONDE A: marco_teorico.tex ---

% --- DIAPOSITIVA 4A: MARCO TEÓRICO - ESPACIO DE POLINOMIOS ---
\begin{frame}{Marco Teórico: Espacio de Polinomios}
	\begin{block}{Espacio de Polinomios }
		Sea $\mathbb{K}$ un cuerpo (normalmente $\mathbb{R}$ o $\mathbb{C}$).
		El \emph{espacio de polinomios en una variable} se define como
		\[
			\mathbb{P}[z]\;=\;
			\Bigl\{\,p(z)=\sum_{k=0}^{n} a_k\,z^{\,k}\;:\;
			a_k\in\mathbb{K},\;n\in\mathbb{N}\Bigr\}.
		\]
		Para cada $n\in\mathbb{N}$ denotamos
		$
			\mathbb{P}_{n}[z]=\{p\in\mathbb{P}[z]:\deg p\le n\},
		$
		de modo que $\mathbb{P}[z]=\bigcup_{n=0}^{\infty}\mathbb{P}_{n}[z]$.
		El conjunto $\{1,z,z^{2},\dots,z^{n}\}$ forma una base de
		$\mathbb{P}_{n}[z]$, luego $\dim\!\bigl(\mathbb{P}_{n}[z]\bigr)=n+1$.
	\end{block}
\end{frame}

% --- DIAPOSITIVA 4B: MARCO TEÓRICO - PRODUCTO INTERNO GENERAL ---
\begin{frame}{Marco Teórico: Producto Interno General}
	\begin{block}{Producto Interno }
		Sea $\mu$ una medida de Borel positiva con soporte infinito $\operatorname{supp}(\mu)=\{z\;\mu(B_r(z))>0\;\forall r>0\}$. Definimos en $\mathbb{P}[z]$ un \emph{producto interno} cualquiera por
		\[
			\langle p,q\rangle
			\;=\;
			\int p(z)\,\overline{q(z)}\,\mathrm d\mu(z),
			\qquad p,q\in\mathbb{P}[z].
		\]
		Esta forma sesqui-lineal satisface:
		\begin{enumerate}
			\item $\langle p,q\rangle=\overline{\langle q,p\rangle}$,
			\item $\langle \alpha p+\beta r,q\rangle
				      =\alpha\langle p,q\rangle+\beta\langle r,q\rangle$,
			\item $\langle p,p\rangle>0$ si $p\neq0$.
		\end{enumerate}
	\end{block}
\end{frame}

% --- NUEVA DIAPOSITIVA: MARCO TEÓRICO - ESPACIO PREHILBERTIANO Y DE HILBERT ---
\begin{frame}{Marco Teórico: Espacio Prehilbertiano y de Hilbert}
	\footnotesize
	\begin{block}{Espacio Prehilbertiano $(X, \langle \cdot, \cdot \rangle)$}
		Espacio vectorial $X$ con producto interno $\langle \cdot, \cdot \rangle: X \times X \to \mathbb{C}$ (o $\mathbb{R}$) tal que:
		\begin{enumerate}
			\item $\langle x, y \rangle = \overline{\langle y, x \rangle}$ (simetría hermitiana)
			\item $\langle \alpha x + \beta y, z \rangle = \alpha \langle x, z \rangle + \beta \langle y, z \rangle$ (linealidad en la primera variable)
			\item $\langle x, x \rangle \ge 0$ y $\langle x, x \rangle = 0 \iff x = 0$ (definida positiva)
		\end{enumerate}
	\end{block}
	\pause
	\begin{alertblock}{Espacio de Hilbert}
		Un espacio prehilbertiano que además es \textbf{completo} respecto a la norma inducida $\|x\| = \sqrt{\langle x, x \rangle}$ (toda sucesión de Cauchy converge).
	\end{alertblock}
\end{frame}

% --- NUEVA DIAPOSITIVA: MARCO TEÓRICO - GRAM-SCHMIDT Y POLINOMIOS ORTONORMALES ---
\begin{frame}{Marco Teórico: Gram-Schmidt y Polinomios Ortonormales}
	\footnotesize
	\begin{block}{Proceso de Gram-Schmidt}
		Sea $\{v_0,v_1,\dots,v_n\}$ un sistema generador de un subespacio $W\subset V$ con producto interno. Construimos recursivamente:
		\[
			u_0=v_0,\qquad
			u_k= v_k-\sum_{j=0}^{k-1}
			\frac{\langle v_k,u_j\rangle}{\langle u_j,u_j\rangle}\,u_j,
			\;k=1,\dots,n.
		\]
		El conjunto ortonormal $\displaystyle e_k=\frac{u_k}{\|u_k\|}$ satisface $\langle e_i,e_j\rangle=\delta_{ij}$ y $\operatorname{span}\{e_0,\dots,e_n\}=W$.
	\end{block}
	\pause
	\begin{alertblock}{Sucesión de Polinomios Ortonormales}
		Dado un producto interno $\langle\cdot,\cdot\rangle$ en $\mathbb{P}[z]$, la \emph{sucesión de polinomios ortonormales} $\{\varphi_n\}_{n\ge 0}$ se obtiene aplicando Gram-Schmidt a la base monomial $\{1,z,z^2,\dots\}$:
		\[
			\deg \varphi_n=n,\qquad
			\langle \varphi_n,\varphi_k\rangle=\delta_{nk}.
		\]
		Cada $\varphi_n$ es único salvo su signo si se exige que su coeficiente principal sea positivo.
	\end{alertblock}
\end{frame}

% --- NUEVA DIAPOSITIVA: MARCO TEÓRICO - CEROS DE POLINOMIOS ---
\begin{frame}{Marco Teórico: Ceros de Polinomios}
	\footnotesize
	\begin{block}{Ceros de un Polinomio }
		Sea $p\in\mathbb{P}[z]\setminus\{0\}$. El \emph{conjunto de ceros} de $p$ es:
		\[
			Z(p)=\bigl\{\,z\in\mathbb{C}\;:\;p(z)=0\bigr\}.
		\]
		Para la sucesión de polinomios ortonormales $\{\varphi_n\}_{n\ge0}$, el \emph{conjunto total de ceros} $Z$ es:
		\[
			Z \;=\; \bigcup_{n=0}^{\infty} Z(\varphi_n).
		\]
	\end{block}
	\pause
	\begin{alertblock}{Importancia}
		La localización y propiedades de estos conjuntos de ceros son un tema central en la teoría de polinomios ortogonales. Para facilitar su estudio, a menudo se recurre a representaciones matriciales.
	\end{alertblock}
\end{frame}

% --- NUEVA DIAPOSITIVA: MARCO TEÓRICO - MATRIZ HERMITIANA Y DEFINIDA POSITIVA ---
\begin{frame}{Marco Teórico: Matriz Hermitiana y Definida Positiva}
	\begin{block}{Matriz Hermitiana \cite{Horn_Johnson_2012}}
		$A \in \mathcal{M}_{n}(\mathbb{C})$ es \textbf{hermitiana} si $A = A^*$ (conjugada transpuesta).
		Si $A \in \mathcal{M}_{n}(\mathbb{R})$, es \textbf{simétrica} si $A = A^t$.
	\end{block}
	\pause
	\begin{alertblock}{Definición Positiva / Semidefinida Positiva}
		Una matriz hermitiana $A$ es:
		\begin{itemize}
			\item \textbf{Definida positiva (HPD):} $xAx^* > 0$ para todo $x \in \mathbb{C}^n, x \neq 0$. (Todos los autovalores $>0$).
			\item \textbf{Semidefinida positiva (HPSD):} $xAx^* \ge 0$ para todo $x \in \mathbb{C}^n$. (Todos los autovalores $\ge 0$).
		\end{itemize}
		$\langle x, y \rangle_A = xAy^*$ define un producto interno si y solo si $A$ es HPD.
	\end{alertblock}
\end{frame}

% --- NUEVA DIAPOSITIVA: MARCO TEÓRICO - MATRIZ DE MOMENTOS ---
\begin{frame}{Marco Teórico: Matriz de Momentos}
	\begin{block}{Matriz de Momentos} % (Definición \ref{def:moment_matrix})
		Sea $\langle\cdot,\cdot\rangle$ un producto interno en $\mathbb{P}[z]$ y $\mathcal{B}=\{z^k\}_{k\ge 0}$ la base monomial.
		La \textbf{matriz de momentos} asociada es $M=(c_{i,j})_{i,j\ge 0}$, donde:
		\[ c_{i,j} := \langle z^i, z^j \rangle, \quad i,j\in\mathbb{N}_0. \]
		Si $\langle p,q \rangle = \int p(z)\overline{q(z)}d\mu(z)$, entonces:
		\[ c_{i,j} = \int z^i \overline{z^j} d\mu(z). \]
		Esta matriz encapsula la información del producto interno respecto a la base monomial.
	\end{block}
\end{frame}

% --- NUEVA DIAPOSITIVA: MARCO TEÓRICO - PRODUCTO INTERNO INDUCIDO POR MATRIZ HPD ---
\begin{frame}{Marco Teórico: Producto Interno Inducido por Matriz HPD}
	\footnotesize
	\begin{block}{Producto Interno Inducido por una Matriz HPD}
		Desde una perspectiva matricial, que facilita el análisis algebraico, una matriz infinita $M=(c_{ij})_{i,j \ge 0}$ Hermitiana Definida Positiva (HPD) induce en $\mathbb{P}[z]$ el siguiente producto interno:
		\begin{equation} \label{eq:producto_matricial}
			\langle p(z), q(z) \rangle_M = \mathbf{v} M \mathbf{w}^{*}
		\end{equation}
		donde:
		\begin{itemize}
			\item $p(z) = \sum_{k=0}^n v_k z^k$ y $q(z) = \sum_{k=0}^m w_k z^k$ son polinomios.
			\item $\mathbf{v} = (v_0, v_1, \ldots, v_n, 0, \ldots)$ y $\mathbf{w} = (w_0, w_1, \ldots, w_m, 0, \ldots)$ son los vectores de coeficientes de $p$ y $q$ en la base monomial. Estos vectores pertenecen a $c_{00}$, el espacio de sucesiones complejas con un número finito de términos no nulos:
			      $c_{00} = \{ (x_k)_{k=0}^{\infty} \in \mathbb{C}^{\mathbb{N}_0} : \exists N \in \mathbb{N}_0 \text{ tal que } x_k = 0 \text{ para todo } k \ge N \}$.
		\end{itemize}
		Los elementos $c_{ij}$ de la matriz $M$ son, consistentemente, los momentos respecto a este producto: $c_{ij} = \langle z^i, z^j \rangle_M$.
		Este producto interno induce la norma Hilbertiana:
		\[ \|p(z)\|^{2}_{M} = \langle p(z), p(z)\rangle_{M}. \]
	\end{block}
\end{frame}

% --- NUEVA DIAPOSITIVA: MARCO TEÓRICO - SECCIÓN PRINCIPAL DE UNA MATRIZ ---
\begin{frame}{Marco Teórico: Sección Principal de una Matriz}
	\footnotesize
	\begin{block}{Sección Principal de Orden $n$}
		Sea $A=(a_{ij})_{i,j\ge 0}$ una matriz infinita. Su \textbf{sección principal de orden $n$} (o truncación) es:
		\[
			A_n = (a_{ij})_{0\le i,j\le n-1} =
			\begin{pmatrix}
				a_{00}    & \dots  & a_{0,n-1}   \\
				\vdots    & \ddots & \vdots      \\
				a_{n-1,0} & \dots  & a_{n-1,n-1}
			\end{pmatrix}.
		\]
		Estas secciones permiten “aproximar” $A$ por operadores de rango finito y realizar cálculos numéricos. Son cruciales para el análisis numérico.
	\end{block}
\end{frame}

% --- CORRESPONDE A: marco_teorico.tex ---
% --- DIAPOSITIVA 5: MARCO TEÓRICO - OPERADOR MULTIPLICACIÓN (I) ---
\begin{frame}{Marco Teórico: Operador Multiplicación $\D(p)=zp$ (I)}
	\begin{block}{Definición y Acotación} \small
		\begin{itemize}
			\item Operador lineal $\D: \mathbb{P}[z] \to \mathbb{P}[z]$ dado por $\D(p(z)) = z \cdot p(z)$.
			\item $\D$ es \textbf{acotado} si existe $C>0$ tal que $\|\D(p)\| \le C \|p\|$ para todo $p \in \mathbb{P}[z]$.
			\item La menor $C$ es la norma del operador $\|\D\|$.
		\end{itemize}
	\end{block}
	\pause
	\begin{alertblock}{Representación Matricial $\mathbf{D}$} \small
		\begin{itemize}
			\item En la base de polinomios ortonormales $\{\varphi_n\}$, $z\varphi_j(z) = \sum_{i=0}^{j+1} d_{i,j} \varphi_i(z)$.
			\item Esto define una matriz infinita $\mathbf{D} = (d_{i,j})_{i,j \ge 0}$.
			\item $\mathbf{D}$ es una matriz de \textbf{Hessenberg superior}: $d_{i,j}=0$ si $i > j+1$.
			      \[
				      \mathbf{D}=\begin{pmatrix}
					      d_{0,0} & d_{0,1} & d_{0,2} & \cdots \\
					      d_{1,0} & d_{1,1} & d_{1,2} & \cdots \\
					      0       & d_{2,1} & d_{2,2} & \cdots \\
					      \vdots  & \vdots  & \ddots  & \ddots
				      \end{pmatrix}
			      \]
		\end{itemize}
	\end{alertblock}
\end{frame}

% --- CORRESPONDE A: marco_teorico.tex ---
% --- DIAPOSITIVA 6: MARCO TEÓRICO - OPERADOR MULTIPLICACIÓN (II) - CONEXIÓN ESPECTRAL ---
\begin{frame}{Marco Teórico: Operador Multiplicación (II) - Conexión Espectral}
	\begin{block}{Propiedades de la Matriz $\mathbf{D}$}
		\begin{itemize}
			\item Si el producto interno procede de una medida real en un intervalo (caso clásico), la relación de recurrencia de tres términos hace que $\mathbf{D}$ sea \textbf{tridiagonal y simétrica} (matriz de Jacobi).
			\item Las secciones principales $\mathbf{D}_n$ son matrices $(n+1)\times(n+1)$ cuya transpuesta tiene como autovalores los ceros de $\varphi_{n+1}$.
			\item El determinante $\det(\mathbf{D}_n - z\mathbf{I}_{n+1})$ es un polinomio proporcional a $\varphi_{n+1}(z)$ (tienen las mismas raíces).
		\end{itemize}
	\end{block}
\end{frame}

% --- CORRESPONDE A: marco_teorico.tex ---
% --- DIAPOSITIVA 6B: MARCO TEÓRICO - OPERADOR MULTIPLICACIÓN (II) - CONEXIÓN ESPECTRAL (DEMOSTRACIÓN) ---
\begin{frame}{Marco Teórico: Operador Multiplicación (II) - Conexión Espectral (Demostración)}
	\begin{block}{Demostración (Idea)} \small
		Si $\varphi_{n+1}(z_0)=0$:
		\begin{align*} z_0\varphi_k(z_0) &= \sum_{i=0}^{k+1} d_{i,k} \varphi_i(z_0) \quad \text{para } k < n \\ z_0\varphi_n(z_0) &= \sum_{i=0}^{n} d_{i,n} \varphi_i(z_0) \quad (\text{ya que } d_{n+1,n}\varphi_{n+1}(z_0)=0) \end{align*}
		Esto se escribe como $z_0 \vec{\Phi}(z_0) = \mathbf{D}_n^t \vec{\Phi}(z_0)$, donde $\vec{\Phi}(z_0) = (\varphi_0(z_0), \dots, \varphi_n(z_0))^t$.
	\end{block}
	\pause \small
	\begin{itemize}
		\item En el caso clásico (medida real), $\mathbf{D}$ es tridiagonal y simétrica (matriz de Jacobi).
		\item Esta conexión permite usar herramientas de álgebra lineal para estudiar los ceros.
	\end{itemize}
\end{frame}

% --- CORRESPONDE A: marco_teorico.tex ---
% --- DIAPOSITIVA 7: MARCO TEÓRICO - PARÁMETROS ESPECTRALES CLAVE ---
\begin{frame}{Marco Teórico: Parámetros Espectrales Clave}\footnotesize
	\begin{block}{Definición: Radio Espectral \cite{dunford1963linear}}
		Si $\mathbf{A}$ es cuadrada, su \textbf{radio espectral} $\rho(\mathbf{A})$ es el máximo de los módulos de sus autovalores:
		$$ \rho(\mathbf{A}) = \max \{ |\lambda| : \lambda \text{ es autovalor de } \mathbf{A} \}. $$
		Para $\mathbf{D}_n$, $\rho(\mathbf{D}_n)$ nos da la localización del cero de $\varphi_{n+1}$ más alejado del origen.
	\end{block}
	\pause
	\begin{block}{Definición: Valores Singulares \cite{Horn_Johnson_2012}}
		Los \textbf{valores singulares} de $\mathbf{A} \in \mathcal{M}_{m \times n}(\mathbb{C})$, denotados $\sigma_k(\mathbf{A})$, se definen a partir de los autovalores $\lambda_k(\mathbf{A}^*\mathbf{A})$ de la matriz HPSD $\mathbf{A}^*\mathbf{A}$ (o $\mathbf{A}^t\mathbf{A}$ si $\mathbf{A}$ es real):
		$$ \sigma_k(\mathbf{A}) = \sqrt{\lambda_k(\mathbf{A}^*\mathbf{A})}. $$
		Se ordenan de forma decreciente:
		$$ \sigma_1(\mathbf{A}) \ge \sigma_2(\mathbf{A}) \ge \dots \ge \sigma_{\min(m,n)}(\mathbf{A}) \ge 0. $$
	\end{block}
\end{frame}

\begin{frame}
	\begin{block}{Proposición: Norma Espectral}
		El \textbf{mayor valor singular} $\sigma_1(\mathbf{A})$ (Criterio de Rayleigh):
		$$
			\sigma_{1}(\mathbf{A})^2 =
			\max \{ \langle\mathbf{A}^{*}\mathbf{A}x,x\rangle: \| x \|_2 =1 \}
			= \max \{\| \mathbf{A} x \|_2^2 : \| x \|_2 =1\}=
			\| \mathbf{A} \|_2^2
		$$
		Para $\mathbf{D}_n$, $\|\mathbf{D}_n\|_2 = \sigma_{1,n}$ mide la ``acotación'' de la sección del operador.
	\end{block}
	\begin{block}{Proposición: Segundo Valor Singular}
		El \textbf{segundo mayor valor singular} $\sigma_2(\mathbf{A})$ (Principio min-max):
		$$
			\sigma_{2}(\mathbf{A})^2 = \min_{\substack{V \subset \mathbb{C}^n \\ \text{codim}(V)=1}} \max \{\| \mathbf{A}x \|_2^2 : x\in V, \| x \|_2 =1\}= \min_{\substack{V \subset \mathbb{C}^n \\ \text{codim}(V)=1}} \| \mathbf{A}_{|V} \|_2^2
		$$
		donde $\mathbf{A}_{|V}$ es la restricción de $\mathbf{A}$ al subespacio $V$. Este valor jugará un papel crucial en nuestro análisis.
	\end{block}
\end{frame}

% --- NUEVA DIAPOSITIVA: MARCO TEÓRICO - NOTACIÓN PARA SECCIONES MATRICIALES ---
\begin{frame}{Marco Teórico: Notación para Secciones Matriciales}
	\begin{block}{Definición: Notación Específica para $\mathbf{D}_n$}
		Sea $\mathbf{D}_{n}$ la sección de la matriz infinita $\mathbf{D}$ que representa al operador $\mathcal{D}$ (multiplicación por $z$) en una base de polinomios ortogonales. Definimos:
		\begin{enumerate}
			\item Valores singulares de $\mathbf{D}_{n}$: $\sigma_{1,n} \ge \sigma_{2,n} \ge \ldots \ge \sigma_{n+1,n}$ (ya que $\mathbf{D}_n$ es $(n+1) \times (n+1)$).
			\item Radio espectral: $\rho(\mathbf{D}_{n})$.
			\item Conjunto de ceros de $\varphi_{n+1}$: $Z_{n+1}$.
		\end{enumerate}
		Esta notación simplificará la discusión sobre el comportamiento asintótico.
	\end{block}
\end{frame}

% --- NUEVA DIAPOSITIVA: MARCO TEÓRICO - LÍMITE DE NORMAS DE SECCIONES ---
\begin{frame}{Marco Teórico: Límite de Normas de Secciones}
	\begin{block}{Proposición: Convergencia de la Norma de las Secciones}
		Sea $\mathcal{D}$ el operador multiplicación por $z$ en $\ell_2$ (espacio de sucesiones de cuadrado sumable) y $\mathbf{D}=(d_{i,j})_{i,j\ge0}$ su matriz infinita. Sea $\mathbf{D}_n=(d_{i,j})_{0\le i,j\le n}$ la sección de orden $n+1$. Entonces:
		\[
			\lim_{n\to\infty}\|\mathbf{D}_n\|_2 \;=\;\|\mathcal{D}\|.
		\]
	\end{block}
	Esta proposición establece una conexión importante entre las normas de las secciones finitas y la norma del operador infinito.
	\pause
	\begin{alertblock}{Consecuencia Directa}
		Utilizando este resultado, y dado que $\|\mathbf{D}_n\|_2 = \sigma_{1,n}(\mathbf{D}_n)$:
		$$
			\lim_{n\to \infty} \sigma_{1,n}(\mathbf{D}_n) = \lim_{n \to \infty} \|\mathbf{D}_n\|_2 = \|\mathcal{D}\|.
		$$
		Este límite siempre existe, pudiendo ser $\infty$ si el operador $\mathcal{D}$ no es acotado.
	\end{alertblock}
\end{frame}

% --- NUEVA DIAPOSITIVA: MARCO TEÓRICO - MEDIDA DE LEBESGUE EN CIRCUNFERENCIA ---
\begin{frame}{Marco Teórico: Medida de Lebesgue en Circunferencia}
	\begin{block}{Definición}
		Sea $S_{r}(a) = \{z\in\mathbb{C} : |z-a|=r\}$ la circunferencia de centro $a$ y radio $r$.
		La \textbf{medida de Lebesgue normalizada} $dm_{a;r}$ en $S_{r}(a)$ se define para $f:S_{r}(a)\to\mathbb{C}$ continua como:
		$$ \int_{S_{r}(a)} f(z)\,dm_{a;r}(z) := \frac{1}{2\pi}\int_{0}^{2\pi} f(a+re^{i\theta})\,d\theta $$
		Así, $dm_{a;r}(S_{r}(a))=1$.
	\end{block}
	\begin{block}{Definición}
		Sea $b \in \mathbb{C}$ y $r > 0$. El \textbf{disco abierto} de centro $b$ y radio $r$ se define como:
		$$ \mathbb{D}_r(b) = \{z \in \mathbb{C} : |z - b| < r\} $$
	\end{block}
\end{frame}

% --- CORRESPONDE A: marco_teorico.tex ---
% --- DIAPOSITIVA 8: MARCO TEÓRICO - PRODUCTOS INTERNOS DE SOBOLEV (I) - DEFINICIÓN ---
\begin{frame}{Marco Teórico: Productos Internos de Sobolev (I) - Definición}
	\begin{alertblock}{Definición General \cite{escribano2025boundedness}}
		Dados $\boldsymbol{\mu}=(\mu_1, \mu_2)$ medidas de Borel positivas (soporte de $\mu_1$ infinito):
		$$ \langle p,q \rangle_S = \int_{\mathbb{C}} p(z)\overline{q(z)}d\mu_1(z) + \int_{\mathbb{C}} p'(z)\overline{q'(z)}d\mu_2(z) $$
	\end{alertblock}
\end{frame}

% --- CORRESPONDE A: marco_teorico.tex ---
% --- DIAPOSITIVA 8B: MARCO TEÓRICO - PRODUCTOS INTERNOS DE SOBOLEV (I) - EJEMPLOS ---
\begin{frame}{Marco Teórico: Productos Internos de Sobolev (I) - Casos}

	\begin{enumerate}
		\item \textbf{Continuo Real:} $\int_a^b p(t)q(t)w_1(t)dt + \int_c^d p'(t)q'(t)w_2(t)dt$
		\item \textbf{Discreto Real:} $\int_a^b p(t)q(t)w_1(t)dt + p'(c)q'(c)$
		\item \textbf{Continuo Complejo (Circunferencias):}
		      $$ \int_{S_{r_b}(b)} p\bar{q}\,dm_{b;r_b} + \int_{S_{r_a}(a)} p'\bar{q}'\,dm_{a;r_a} $$
		      donde $dm_{c;r}$ es la medida de Lebesgue normalizada en $S_r(c)$.
		\item \textbf{Discreto Complejo (Circunferencia y Átomo):}
		      $$ \int_{S_{r_b}(b)} p\bar{q}\,dm_{b;r_b} + p'(c)\overline{q'(c)} $$
	\end{enumerate}
	\pause \small
	\begin{alertblock}{Términos Discretos y Medida de Dirac}
		Para términos discretos, $\mu_2$ es una \textbf{medida de Dirac} $\delta_c$.
	\end{alertblock}
\end{frame}

% --- NUEVA DIAPOSITIVA: MARCO TEÓRICO - MEDIDA DE DIRAC ---
\begin{frame}{Marco Teórico: Medida de Dirac}
	\begin{block}{Definición}
		Sea $a\in\mathbb{C}$. La \textbf{medida de Dirac} en $a$, denotada $\delta_{a}$, es la medida de Borel:
		\[ \delta_{a}(A) = \begin{cases} 1, & a\in A \\ 0, & a\notin A \end{cases} \quad (A\subset\mathbb C \text{ Borel}) \]
		Para $f:\mathbb{C}\to\mathbb{C}$ medible: $\displaystyle \int_{\mathbb C} f(z)\,d\delta_{a}(z)=f(a)$.
	\end{block}
	\pause
	\begin{alertblock}{Aplicación a Polinomios}
		En particular, $|p(a)|^{2} = \int_{\mathbb C} |p(z)|^{2}\,d\delta_{a}(z)$.
		Para ponderar, se usa $\lambda\,\delta_{a}\; (\lambda>0)$, con $\int f\,d(\lambda\delta_{a})=\lambda\,f(a)$.
	\end{alertblock}
\end{frame}

% --- CORRESPONDE A: marco_teorico.tex ---
% --- DIAPOSITIVA 9B: MARCO TEÓRICO - PRODUCTOS INTERNOS DE SOBOLEV (IIB) ---
\begin{frame}{Marco Teórico: Productos Sobolev (IIb) - Matriz $\boldsymbol{\Lambda}$ y Casos Discretos}
	\begin{block}{Matriz de Derivación $\boldsymbol{\Lambda}$} \small
		$\boldsymbol{\Lambda}$ es la matriz del operador derivación $p(z) \mapsto p'(z)$ en la base monomial $\{z^k\}$:
		\[ (z^k)' = k z^{k-1} \]
		\[
			\boldsymbol{\Lambda} = \begin{pmatrix}
				0      & 0      & 0      & \cdots \\
				1      & 0      & 0      & \cdots \\
				0      & 2      & 0      & \cdots \\
				\vdots & \vdots & \ddots & \ddots
			\end{pmatrix}
		\]
	\end{block}
	\pause
	\begin{alertblock}{Caso Discreto con Medida de Dirac} \small
		\begin{itemize}
			\item Si $\mu_2 = \delta_a$, su matriz de momentos $M(a)$ tiene entradas $(a^n \bar{a}^m)$.
		\end{itemize}
	\end{alertblock}
\end{frame}

% --- NUEVA DIAPOSITIVA: MARCO TEÓRICO - MATRIZ DE MOMENTOS DE DIRAC ---
\begin{frame}{Marco Teórico: Matriz de Momentos de Dirac $M(a)$}
	\begin{block}{Definición \cite{escribano2025boundedness}}
		Sea $\delta_{a}$ la medida de Dirac en $a\in\mathbb C$. Sus momentos son $c_{n,m} = a^{\,n}\,\overline{a}^{\,m}$.
		La \textbf{matriz de momentos} $M(a)=M(\delta_{a})$ es:
		\(
		M(a) = \begin{pmatrix}
			1      & \overline{a}      & \overline{a}^{2}  & \dots  \\
			a      & |a|^{2}           & a\overline{a}^{2} & \dots  \\
			a^{2}  & a^{2}\overline{a} & |a|^{4}           & \dots  \\
			\vdots & \vdots            & \vdots            & \ddots
		\end{pmatrix}
		\)
		$M(a)$ es Hermitiana y semidefinida positiva (HPSD), no HPD si $a \neq 0$. No define un producto interno por sí sola.
	\end{block}
\end{frame}

% --- NUEVA DIAPOSITIVA: MARCO TEÓRICO - RADIO DEL SOPORTE COMBINADO R ---
\begin{frame}{Marco Teórico: Radio del Soporte Combinado $R$}
	\begin{block}{Definición del Radio $R$}
		En un espacio de Sobolev asociado a dos medidas se define el número $R$:
		$$R = \max \{ |z| : z \in \operatorname{supp}(\mu_1) \cup \operatorname{supp}(\mu_2) \}$$
	\end{block}
	\pause
	\begin{block}{Casos Particulares} \small
		\begin{enumerate}
			\item \textbf{Caso continuo real:} $R = \max \{ |a|, |b|, |c|, |d| \}$
			\item \textbf{Caso discreto real:} $R = \max \{|a|, |b|, |c|\}$
			\item \textbf{Caso continuo complejo:} $R = \max \{ |a| + r_a, |b| + r_b \}$
			\item \textbf{Caso discreto complejo:} $R = \max \{ |a| + r_a, |c| \}$
		\end{enumerate}
	\end{block}
	\pause
	\begin{alertblock}{Propiedad Fundamental} \small
		Para productos internos asociados a una sola medida: $R = \max \{ |z| : z \in \operatorname{supp}(\mu) \}$.
		Cuando la medida tiene soporte acotado: $\lim_{n\to \infty} \sigma_{1,n} = \|\mathcal{D}\| = R$.
	\end{alertblock}
\end{frame}

% --- NUEVA DIAPOSITIVA: MARCO TEÓRICO - PRODUCTO INTERNO DE SOBOLEV MATRICIAL ---
\begin{frame}{Marco Teórico: Producto Interno de Sobolev Matricial}
	\begin{block}{Definición Generalizada \cite{escribano2025boundedness}}
		Sean $\{\mathbf{M}_1, \mathbf{M}_2\}$ matrices HPSD, con $\mathbf{M}_1$ HPD.
		El \textbf{producto interno de Sobolev matricial} $\langle \cdot, \cdot \rangle_{\mathbf{M}_{S}}$ se define como:
		$$ \langle p, q \rangle_{\mathbf{M}_{S}} = \langle p, q \rangle_{\mathbf{M}_1} + \langle p', q' \rangle_{\mathbf{M}_2} $$
		donde $\langle u, v \rangle_{\mathbf{M}_k} = \vec{u}^t \mathbf{M}_k \overline{\vec{v}}$.
		Equivalentemente, la matriz de momentos de Sobolev es:
		$$ \mathbf{M}_S = \mathbf{M}_1 + \boldsymbol{\Lambda}\,\mathbf{M}_2\, \boldsymbol{\Lambda}^* $$
	\end{block}
\end{frame}

% --- NUEVA DIAPOSITIVA: MARCO TEÓRICO - OBSERVACIÓN PRODUCTO MATRICIAL ---
\begin{frame}{Marco Teórico: Observación Producto Matricial}
	\begin{block}{Casos Particulares}
		La definición de producto interno de Sobolev matricial:
		$$ \langle p, q \rangle_{\mathbf{M}_{S}} = \langle p, q \rangle_{\mathbf{M}_1} + \langle p', q' \rangle_{\mathbf{M}_2} $$
		recupera los ejemplos comunes (reales, complejos, continuos, discretos) cuando $\mathbf{M}_1$ y $\mathbf{M}_2$ son matrices de momentos de medidas $\mu_1$ y $\mu_2$ respectivamente.
		\begin{itemize}
			\item $\mathbf{M}_1 = M(\mu_1)$
			\item $\mathbf{M}_2 = M(\mu_2)$ (puede ser $M(\delta_a)$ para términos discretos)
		\end{itemize}
	\end{block}
\end{frame}

% --- CORRESPONDE A: marco_teorico.tex ---
% --- DIAPOSITIVA 10: MARCO TEÓRICO - PROBLEMA CENTRAL Y METODOLOGÍA ---
\begin{frame}{Marco Teórico: Problema Central y Metodología}
	\begin{exampleblock}{Problema Central de Investigación} % Reducido tamaño
		Sean $\mu_1, \mu_2$ medidas con soporte compacto. Sea $Z$ el conjunto de ceros de los polinomios ortonormales de Sobolev $\{\varphi_n\}$ asociados.
		\vspace{0.1cm}
		\centering \textbf{¿Es $Z$ acotado?}
		\vspace{0.1cm}
		Esta pregunta no está completamente resuelta, especialmente si los soportes son disjuntos o no hay "dominación secuencial".
	\end{exampleblock}
\end{frame}


% --- CORRESPONDE A: marco_teorico.tex ---
% --- DIAPOSITIVA 10B: MARCO TEÓRICO - OBJETIVO Y METODOLOGÍA ---
\begin{frame}{Marco Teórico: Objetivo Específico y Metodología}
	\begin{block}{Objetivo Específico del Marco Teórico} % Reducido tamaño
		Determinar condiciones para la acotación de $Z$ relacionando:
		\begin{itemize}
			\item Acotación de $Z$ (i.e., $\sup_n \rho(\mathbf{D}_n) < \infty$).
			\item Acotación del operador $\mathcal{D}$ (i.e., $\lim_n \sigma_{1,n} = \|\mathcal{D}\| < \infty$).
			\item Comportamiento asintótico de $\rho(\mathbf{D}_n)$, $\sigma_{1,n}$ y $\sigma_{2,n}$.
		\end{itemize}
	\end{block}
	\pause
	\begin{block}{Metodología Computacional} % Reducido tamaño
		Usar \texttt{Maple} para:
		\begin{enumerate}
			\item Generar matrices $\mathbf{D}_n$.
			\item Calcular $\rho(\mathbf{D}_n)$, $\sigma_{1,n}$, $\sigma_{2,n}$ y analizar su convergencia.
			\item Visualizar ceros y contrastar con cotas teóricas.
		\end{enumerate}
	\end{block}
\end{frame}

% --- CORRESPONDE A: marco_teorico.tex ---
% --- DIAPOSITIVA 11A.1: MARCO TEÓRICO - ESTADO DEL ARTE (IA) - ACOTACIÓN DE CEROS ---
\begin{frame}{Marco Teórico: Estado del Arte (Ia) - Acotación de Ceros}
	\begin{block}{Proposición 4 \cite{escribano2025boundedness}}
		Sea \textbf{M} una matriz HPD y $\{\varphi_n(z)\}_{n=0}^{\infty}$ la secuencia de polinomios ortonormales asociados al producto interno inducido por \textbf{M}. Si $\mathcal{D}$ está acotado en $\mathbb{P}[z]$, el conjunto de ceros de los polinomios ortonormales asociados a \textbf{M} está acotado. Además,
		$$Z=\{z\in\mathbb{C}: \exists n \in \mathbb{N},\varphi_n(z)=0\} \subset \{z\in\mathbb{C}:|z|<\|\mathcal{D}\|\}.$$
	\end{block}
\end{frame}

% --- NUEVA DIAPOSITIVA: DEMOSTRACIÓN PROPOSICIÓN 4 ---
\begin{frame}{Marco Teórico: Demostración Proposición 4}
	\begin{proof} \footnotesize
		Asumamos que para $z_0 \in Z$ existe un $n \in \mathbb{N}$ tal que $\varphi_n(z_0)=0$. Luego, $\varphi_n(z)=(z-z_0)P_{n-1}(z)$, $0 \neq P_{n-1} \in \mathbb{P}_{n-1}[z]$.

		Por lo tanto, $\varphi_n(z)=zP_{n-1}(z)-z_0P_{n-1}(z)$ y porque $\varphi_n(z)$ y $P_{n-1}(z)$ son ortogonales en $\textbf{M}$ se sigue que:
		$$\|zP_{n-1}(z)\|^2=\langle\varphi_n(z)+z_0P_{n-1}(z),\varphi_n(z)+z_0P_{n-1}(z)\rangle$$
		$$= \|\varphi_n(z)\|^2+|z_0|^2\|P_{n-1}(z)\|^2$$

		Por lo que:
		$$|z_0|^2\|P_{n-1}(z)\|^2=\|zP_{n-1}(z)\|^2- \|\varphi_n(z)\|^2<\|zP_{n-1}(z)\|^2$$
		$$=\|\mathcal{D}(P_{n-1})\|^2\leq\|\mathcal{D}\|^2\|P_{n-1}(z)\|^2$$

		Como $\|P_{n-1}(z)\| \neq 0$, se concluye que $|z_0|^2<\|\mathcal{D}\|$.
	\end{proof}
\end{frame}

% --- CORRESPONDE A: marco_teorico.tex ---
% --- DIAPOSITIVA 11A.2: MARCO TEÓRICO - ESTADO DEL ARTE (IA) - DOMINACIÓN SECUENCIAL ---
\begin{frame}{Marco Teórico: Estado del Arte (Ia) - Dominación Secuencial (Definición y Corolario)}
	\begin{block}{Definición 22 \cite{lagomasino1999zero}}
		Diremos que el producto de Sobolev con dos medidas es secuencialmente dominado si
		$$\operatorname{supp}(\mu_2) \subset \operatorname{supp}(\mu_{1})$$
		y
		$$d\mu_2 = f_{1}d\mu_{1}, \qquad f_{1} \in L_{\infty}(\mu_{1}).$$
	\end{block}
	\pause
	\begin{block}{Corolario 1} \footnotesize
		El producto interno de Sobolev
		$$ \langle p(t), q(t) \rangle = \int_{a}^{b} p(t) q(t) w_1(t) \, dt + \int_{c}^{d} p'(t) q'(t) w_2(t) \, dt $$
		es secuencialmente dominado cuando se cumple que $[c,d] \subset [a,b]$ y $\exists c > 0, w_2(t)\leq c \cdot w_1(t)$ continuos.
	\end{block}
\end{frame}

% --- NUEVA DIAPOSITIVA: MARCO TEÓRICO - DEFINICIÓN 23 MATRICES SECUENCIALMENTE DOMINADAS ---
\begin{frame}{Marco Teórico: Estado del Arte (Ia) - Matrices Secuencialmente Dominadas}
	\begin{block}{Definición 23 \cite{escribano2025boundedness}}
		Sea $\{\textbf{M}_1,\textbf{M}_2\}$ un conjunto de matrices HSPD. Diremos que $\{\textbf{M}_1,\textbf{M}_2\}$ es secuencialmente dominada si existe una constante positiva $C$ tal que
		$$\textbf{M}_2 \leq C \cdot \textbf{M}_{1}.$$
		donde $\textbf{A} \leq \textbf{B}$ significa que $v\textbf{A}v^* \leq v\textbf{B}v^*$ para todo $v \in c_{00}$.
	\end{block}
	\pause
	\begin{alertblock}{Consecuencia} \footnotesize
		Como consecuencia de \cite{escribano2025zeros} Proposición 5, se tiene que el producto Sobolev
		$$\langle p(z) , q(z) \rangle_{S}\;=\;
			\int_{S_{r_b}(b)} p(z)\,\overline{q(z)} \, dm_{b;r_b}
			\;+\;
			\int_{S_{r_a}(a)} p'(z)\,\overline{q'(z)} \, dm_{a;r_a}.$$
		es matricialmente secuencialmente dominado cuando se cumple que $S_{r_a}(a)\subset \mathbb{D}_{r_b}(b)$.
	\end{alertblock}
\end{frame}

% --- NUEVA DIAPOSITIVA: MARCO TEÓRICO - TEOREMA DE OPERADOR ACOTADO Y COROLARIO ---
\begin{frame}{Marco Teórico: Estado del Arte (Ib) - Teorema de Operador Acotado}
	\begin{block}{Teorema \cite{escribano2025boundedness}} \footnotesize
		Sea $\{\textbf{M}_1,\textbf{M}_2\}$ un conjunto de matrices HPD secuencialmente dominado. Asuma que para $j=1,2$ el operador multiplicación $\mathcal{D}_j$ está acotado en $(\mathbb{P}[z],\langle \cdot,\cdot \rangle)_{\textbf{M}_j}$. Considera el producto matricial de Sobolev siguiente:
		$$\langle p(t),q(t) \rangle_{\textbf{M}_S}=\langle p(t),q(t) \rangle_{\textbf{M}_1}+\langle p'(t),q'(t) \rangle_{\textbf{M}_2}.$$
		Entonces, el operador multiplicación $\mathcal{D}$ está acotado en $(\mathbb{P}[z],\langle \cdot,\cdot \rangle)_{\textbf{M}_S}.$
	\end{block}
	\pause
	\begin{alertblock}{Corolario} \footnotesize
		Para el producto interno de Sobolev
		$$\langle p(z) , q(z) \rangle_{S}\;=\;
			\int_{S_{r_b}(b)} p(z)\,\overline{q(z)} \, dm_{b;r_b}
			\;+\;
			\int_{S_{r_a}(a)} p'(z)\,\overline{q'(z)} \, dm_{a;r_a}.$$
		el operador multiplicación $\mathcal{D}$ está acotado cuando $S_{r_a}(a)\subset \mathbb{D}_{r_b}(b)$. Por tanto el conjunto $Z$ está acotado.
	\end{alertblock}
\end{frame}

% --- NUEVA DIAPOSITIVA: MARCO TEÓRICO - EXTENSIÓN SIN DOMINACIÓN SECUENCIAL ---
\begin{frame}{Marco Teórico: Estado del Arte (Ib) - Extensión sin Dominación Secuencial}
	\begin{alertblock}{Resultado Adicional \cite{escribano2025zeros}} \scriptsize
		Para el producto Sobolev
		$$
			\langle p(z) , q(z) \rangle_{S}\;=\;
			\int_{S_{r_b}(b)} p(z)\,\overline{q(z)} \, dm_{b;r_b}
			\;+\;
			\int_{S_{r_a}(a)} p'(z)\,\overline{q'(z)} \, dm_{a;r_a}.
		$$
		El operador multiplicación $\mathcal{D}$ está acotado cuando se cumple que $a \in \mathbb{D}_{r_b}(b)$ y $r_a > 0$, \textbf{incluso cuando no estamos en un caso matricialmente secuencialmente dominado}.

		Este resultado muestra que hay casos donde la acotación del operador puede obtenerse bajo condiciones geométricas más generales que la dominación secuencial.
	\end{alertblock}
\end{frame}

% --- CORRESPONDE A: marco_teorico.tex ---
% --- DIAPOSITIVA 12A.1: MARCO TEÓRICO - ESTADO DEL ARTE (IIA) - PUNTOS DE EVALUACIÓN ACOTADOS (DEFINICIÓN) ---
\begin{frame}{Marco Teórico: Estado del Arte (IIa) - Puntos de Evaluación Acotados}
	Para productos con términos discretos como $|p'(a)|^2$.
	\begin{block}{Punto de Evaluación Acotado (bpe) \cite{escribano2025zeros}} \small
		$a \in \mathbb{C}$ es \textbf{bpe} de una medida $\mu$ (o matriz HPD $M$) si $\exists C>0$:
		$$ |p(a)|^2 \le C \int |p(z)|^2 d\mu(z) \quad (\text{o } |p(a)|^2 \le C \|p\|_M^2) \quad \forall p \in \mathbb{P}[z] $$
	\end{block}
\end{frame}

% --- NUEVA DIAPOSITIVA: MARCO TEÓRICO - ENVOLTURA CONVEXA POLINÓMICA ---
\begin{frame}{Marco Teórico: Estado del Arte (IIa) - Envoltura Convexa Polinómica}
	\begin{block}{Definición: Envoltura Convexa Polinómica \cite{escribano2025zeros}} \small
		Para un conjunto compacto $K \subset \mathbb{C}$, la \textbf{envoltura convexa polinómica} de $K$, denotada por $P_C(K)$, se define como:
		\[
			P_C(K) = \left\{ z \in \mathbb{C} :\, |p(z)| \leq \max_{\xi\in K}|p(\xi)|\ \text{ para todo } p\in\mathbb{P}[z] \right\}
		\]
	\end{block}
	\pause
	\begin{alertblock}{Propiedades Geométricas \cite{Conway1978, Gamelin1969}} \footnotesize
		\begin{enumerate}
			\item Si $K$ es unión de intervalos reales, $P_C(K) = K$.
			\item Si $K = S_r(c)$ (circunferencia), $P_C(K) = \overline{\mathbb{D}_r(c)}$ (disco cerrado).
		\end{enumerate}
		Estos resultados conectan la geometría del soporte de la medida con las propiedades analíticas de los puntos de evaluación.
	\end{alertblock}
\end{frame}

% --- NUEVA DIAPOSITIVA: MARCO TEÓRICO - LEMA DE PUNTOS DE EVALUACIÓN ACOTADOS ---
\begin{frame}{Marco Teórico: Estado del Arte (IIa) - Lema de bpe}
	\begin{block}{Lema \cite{escribano2025zeros}}
		Todo \textbf{punto de evaluación acotado} de una medida $\mu$ está contenido en la envoltura polinómica convexa del soporte de $\mu$:
		$$ \text{Si } a \text{ es bpe de } \mu, \text{ entonces } a \in P_C(\operatorname{supp}(\mu)) $$
		Además, todo \textbf{punto de evaluación acotado} de una medida $\mu$ es un punto de evaluación de la matriz de momentos asociada a $\mu$.
	\end{block}
\end{frame}

% --- NUEVA DIAPOSITIVA: MARCO TEÓRICO - BPE PARA MATRICES ---
\begin{frame}{Marco Teórico: Estado del Arte (IIa) - bpe para Matrices}
	\begin{block}{Definición: bpe para Matrices \cite{escribano2025zeros}} \small
		Sea $M$ una matriz HPD y $a\in\mathbb{C}$.
		Decimos que $a$ es un \textbf{punto de evaluación acotado} (bpe) de $M$ si existe una constante $C>0$ tal que:
		\[ |p(a)|^{2}\;\le\;C\,\|p(z)\|_{M}^{2}, \quad p(z)\in\mathbb{P}[z] \]
	\end{block}
	\pause
	\begin{alertblock}{Proposición: Equivalencia con Dominacion Secuencial \cite{escribano2025zeros}} \small
		Sea $M$ una matriz HPD y $a\in\mathbb{C}$. Son equivalentes:
		\begin{enumerate}
			\item $a$ es un \textbf{punto de evaluación acotado} de $M$;
			\item el conjunto $\{M,\,M(a)\}$ es secuencialmente dominado.
		\end{enumerate}
		donde $M(a)$ es la matriz de momentos de $\delta_a$.
	\end{alertblock}
\end{frame}

% --- NUEVA DIAPOSITIVA: MARCO TEÓRICO - PROPOSICIÓN DE DERIVADA ACOTADA ---
\begin{frame}{Marco Teórico: Estado del Arte (IIa) - Preservación bajo Derivada}
	\begin{block}{Proposición: Preservación de Acotación \cite{escribano2025zeros}} \footnotesize
		Sea $M$ una matriz HPD tal que el operador de multiplicación $\mathcal{D}$ es acotado en $(\mathbb{P}[z],\|\cdot\|_{M})$, y sea $a\in\mathbb{C}$ punto de evaluación acotado de $M$.

		Entonces $\mathcal{D}$ también es acotado en $\mathbb{P}[z]$ respecto a la norma:
		\[
			\|p(z)\|_{\mathbf{M}_{S}}^{2}\;=\;\|p(z)\|_{M}^{2}\;+\;\bigl|p'(a)\bigr|^{2}
		\]
	\end{block}
\end{frame}

% --- NUEVA DIAPOSITIVA: MARCO TEÓRICO - COROLARIO PARA CIRCUNFERENCIAS ---
\begin{frame}{Marco Teórico: Estado del Arte (IIa) - Corolario para Circunferencias}
	\begin{block}{Corolario: Caso Discreto en Circunferencias \cite{escribano2025zeros}} \footnotesize
		Sea $m_{b;r_b}$ una medida de Lebesgue soportada en una circunferencia $S_{r_b}(b)$.
		Supongamos que $a\in \mathbb{D}_{r_b}(b)$ y consideremos la norma de Sobolev discreta:
		\[
			\|p\|_{S}^{2}= \int_{S_{r_b}(b)} |p(z)|^{2}\,dm_{b;r_b}\;+\;\bigl|p'(a)\bigr|^{2}, \quad p\in\mathbb{P}[z]
		\]
		Entonces el operador multiplicación $\mathcal{D}\colon p(z)\mapsto z\,p(z)$ es acotado en $(\mathbb{P}[z],\|\cdot\|_{S})$.

		\textbf{Consecuencia:} El conjunto de ceros de los polinomios ortogonales de Sobolev asociados está acotado.
	\end{block}
\end{frame}

% --- NUEVA DIAPOSITIVA: MARCO TEÓRICO - MÉTODOS ALTERNATIVOS ---
\begin{frame}{Marco Teórico: Estado del Arte (IIIa) - Métodos Alternativos}
	\begin{block}{Enfoques Alternativos para la Acotación de Ceros} \footnotesize
		Además de las técnicas basadas en matrices de momentos y acotación de operadores, existen otros enfoques:
		\begin{itemize}
			\item \textbf{Enfoque geométrico:} Análisis directo de las propiedades de los soportes de las medidas. Cuando $\operatorname{supp}(\mu_1)$ y $\operatorname{supp}(\mu_2)$ son disjuntos, se obtienen cotas explícitas usando análisis complejo y teoría de potencial.
			\item \textbf{Relaciones de recurrencia:} Propiedades asintóticas de los coeficientes de las relaciones de tres términos de los polinomios ortogonales de Sobolev.
			\item \textbf{Desigualdades integrales:} Técnicas basadas en métodos variacionales para la distribución asintótica de los ceros.
		\end{itemize}
		Estos métodos son útiles para información cualitativa sobre el comportamiento de los ceros.
	\end{block}
\end{frame}

% --- NUEVA DIAPOSITIVA: MARCO TEÓRICO - TEOREMA 2 SOPORTES DISJUNTOS ---
\begin{frame}{Marco Teórico: Estado del Arte (IIIb) - Teorema para Soportes Disjuntos}
	\begin{block}{Teorema 2 (López Lagomasino et al. \cite{lagomasino2001sobolev})} \footnotesize
		En un producto Sobolev del Caso Real Continuo, si $\operatorname{co}(\operatorname{supp}(\mu_1)) \cap \operatorname{co}(\operatorname{supp}(\mu_2)) = \emptyset$, entonces los ceros de $\varphi_n$ están contenidos en el disco centrado en el punto de $\operatorname{co}(\operatorname{supp}(\mu_2))$ más alejado de $\operatorname{co}(\operatorname{supp}(\mu_1))$ y radio el diámetro de $\operatorname{co}(\operatorname{supp}(\mu_1)) \cup \operatorname{co}(\operatorname{supp}(\mu_2))$.

		Donde $\operatorname{co}(K)$ es la envoltura convexa de un conjunto $K$:
		\[
			\operatorname{co}(K) = \left\{\sum_{j=1}^m \lambda_j x_j : x_j\in K, \lambda_j\ge 0, \sum_{j=1}^m \lambda_j =1, m\in \mathbb{N}\right\}
		\]
		es decir, el menor conjunto convexo que contiene a $K$ o, equivalentemente, la intersección de todos los subconjuntos convexos cerrados que contienen a $K$.
	\end{block}
\end{frame}

